\documentclass[11pt]{article}

\usepackage[margin=1in]{geometry}
\usepackage[T1]{fontenc}
\usepackage[utf8]{inputenc}
\usepackage{lmodern}
\usepackage{microtype}
\usepackage{amsmath,amssymb,amsthm,mathtools}
\usepackage{booktabs}
\usepackage{enumitem}
\usepackage[hidelinks]{hyperref}
\usepackage[nameinlink,capitalise]{cleveref}

\newcommand{\C}{\mathbb{C}}
\newcommand{\R}{\mathbb{R}}
\newcommand{\cA}{\mathcal{A}}
\newcommand{\cB}{\mathcal{B}}
\newcommand{\cC}{\mathcal{C}}
\newcommand{\cD}{\mathcal{D}}
\newcommand{\cE}{\mathcal{E}}
\newcommand{\cG}{\mathcal{G}}
\newcommand{\cM}{\mathcal{M}}
\newcommand{\cR}{\mathcal{R}}
\newcommand{\cS}{\mathcal{S}}
\newcommand{\cX}{\mathcal{X}}
\newcommand{\cY}{\mathcal{Y}}
\newcommand{\dd}{\mathrm{d}}
\newcommand{\id}{\mathrm{id}}
\newcommand{\im}{\mathrm{im}}
\newcommand{\HS}{\mathrm{HS}}
\newcommand{\MTT}{\mathrm{MTT}}
\newcommand{\SU}{\mathrm{SU}}
\newcommand{\Tr}{\mathrm{Tr}}
\newcommand{\Z}{\mathbb{Z}}

\numberwithin{equation}{section}

\theoremstyle{definition}
\newtheorem{definition}{Definition}[section]
\newtheorem{assumption}[definition]{Assumption}
\newtheorem{example}[definition]{Example}

\theoremstyle{plain}
\newtheorem{theorem}[definition]{Theorem}
\newtheorem{proposition}[definition]{Proposition}
\newtheorem{corollary}[definition]{Corollary}
\newtheorem{lemma}[definition]{Lemma}

\theoremstyle{remark}
\newtheorem{remark}[definition]{Remark}

\title{Modal Triplet Theory and Heterotic Vacuum Selection:\\
What Finite Ansatz Calculations Prove, What Attraction Adds,\\
and What Physical Selection Still Requires}
\author{Peter Nero}
\date{Version 2, July 2026}

\begin{document}
\maketitle

\begin{abstract}
Heterotic compactification papers often use the word ``selection'' for several
different mathematical achievements: satisfying projected field equations,
finding an isolated point in a finite ansatz, proving attraction under a
specified flow, or assigning the realized vacuum.  These implications are not
equivalent.  We formulate a selection ladder for Modal Triplet Theory (MTT)
and the Hull--Strominger system.  A finite invariant residual proves a full
solution only when the tested residual is complete on the ansatz and every
discarded equation vanishes.  An isolated ansatz solution need not be unique
in the full configuration space.  Attraction requires an actual evolution
and a basin estimate, while physical vacuum selection additionally requires
a preparation law, branch rule, or equivalent source data.  We apply this
distinction to two former case studies.  The diagonal Iwasawa geometry
retains valid local balanced and torsion calculations, but the printed
rank-three bundle, anomaly match, generation count, and normalized Yukawa do
not survive the global bundle audit.  The Lens--Nil model is non-integrable
and therefore is not a Hull--Strominger compactification.  The exact
\(q=79\) arithmetic branch, finite Cech data, and selected rank-two HYM
certificates remain meaningful at their declared tiers, but they do not yet
select a physical visible--hidden compactification.  The result is a
rigorous vocabulary and completion contract, not a claim that the heterotic
landscape has already been uniquely selected.
\end{abstract}

\section*{Revision note: Version 2}
\paragraph{Supersedes.} Version 1, DOI
\href{https://doi.org/10.5281/zenodo.17072927}
{10.5281/zenodo.17072927}.

\paragraph{Reason.}  The former paper defined its fixed-point consistency
condition by the compactification equations themselves and then called their
componentwise solution a selection theorem.  It also inherited invalid
Iwasawa bundle and Yukawa claims and treated a non-integrable Lens--Nil
structure as a heterotic compactification.

\paragraph{Resolution.}  This version separates projected consistency, full
equation closure, ansatz-local isolation, dynamical attraction, and physical
selection.  It proves the exact implication structure, imports the audited
status of the two case studies, retains the current \(q=79\) finite evidence
at its stated tier, and identifies the remaining physical selection
certificate.  Revision history is kept here rather than in the abstract.

\paragraph{Retained content.}  The diagonal Iwasawa balanced geometry and
local torsion calculation, the usefulness of invariant residual tests, and
the exact finite \(q=79\), Cech, and rank-two HYM evidence are retained at
their declared scope.

\paragraph{Remaining boundary.}  A physical compactification-selection claim
still requires the common rank-three visible--hidden Fu--Yau tuple, its
differential Bianchi and gerbe data, the connection-preserving MTT bridge, an
attraction domain, and a source or preparation rule.

\tableofcontents

\section{Why the word selection needs a type}

The vacuum problem is not merely to solve a set of equations.  A field
configuration may solve the equations in a restricted family but fail an
equation omitted by that family.  It may be a genuine solution but be one
among many disconnected branches.  It may be isolated but dynamically
unstable.  It may be attracting but reached only from a small set of initial
conditions.  Finally, a theory may identify all basins without supplying a
law that says which basin is realized.

These distinctions matter especially in heterotic compactification.  The
Hull--Strominger system couples complex geometry, a Hermitian form, a
holomorphic volume form, visible and hidden holomorphic bundles, HYM
connections, torsion, and the differential Green--Schwarz identity on one
common compact complex threefold
\cite{Strominger1986,Hull1986,DelaOssaSvanes2014}.  Solving two scalar
coefficients in a left-invariant truncation is useful, but it is not
automatically a construction of that global tuple.

The first version of this paper obscured this point.  Its ``FCC'' condition
was defined to contain balance, primitivity, quantization, and the Bianchi
identity.  The statement that FCC was equivalent to those same
componentwise equations was therefore a correct identity of definitions, not
a mechanism selecting one physical vacuum.  The purpose of the present
version is to say exactly what each layer does prove.

\section{The lower target: one Hull--Strominger tuple}

Let \(X\) be a compact complex threefold with nowhere-vanishing holomorphic
\((3,0)\)-form \(\Omega\).  Let \(\omega\) be a positive Hermitian form, let
\((V,A)\) and \((W,A_W)\) be visible and hidden holomorphic bundles with HYM
connections, and let \(\nabla\) be the declared tangent-bundle connection.
Suppressing conventional trace normalizations, the target equations include
\begin{align}
\dd(\|\Omega\|_\omega\,\omega^2)&=0, \label{eq:balanced}\\
F_A^{0,2}=F_{A_W}^{0,2}&=0, \qquad
F_A\wedge\omega^2=F_{A_W}\wedge\omega^2=0, \label{eq:hym}\\
H&=\mathrm{i}(\bar\partial-\partial)\omega, \label{eq:torsion}\\
\dd H&=\frac{\alpha'}{4}
\left(\Tr R_\nabla\wedge R_\nabla
-\Tr F_A\wedge F_A-\Tr F_{A_W}\wedge F_{A_W}\right).
\label{eq:bianchi}
\end{align}
Global flux, gerbe, stability, quantization, and worldsheet requirements add
further data.  Established solutions show that this target is mathematically
nonempty, including Fu--Yau-type constructions and stable-bundle
perturbations \cite{FuYau2008,AndreasGarciaFernandez2012,
FinoGrantcharovVezzoni2021}.  Invariant solutions on complex Lie groups also
exist, but their existence does not validate every proposed invariant
bundle \cite{FeiYau2015,OtalUgarteVillacampa2016}.

Write the full typed configuration space as
\[
\cX_{\HS}=\{(X,\Omega,\omega,V,A,W,A_W,\nabla,H,\ldots)\}/\cG,
\]
where \(\cG\) contains the declared gauge and geometric equivalences.  Let
\[
\cR_{\HS}:\cX_{\HS}\longrightarrow\cY_{\HS}
\]
collect all residuals, including \cref{eq:balanced,eq:hym,eq:bianchi} and
the global conditions selected for the problem.  The full solution set is
\[
\cS_{\HS}=\cR_{\HS}^{-1}(0).
\]
This definition is deliberately demanding: all entries refer to one carrier
and one compatible collection of bundles and connections.

\section{The selection ladder}

\subsection{Five distinct levels}

\begin{definition}[Selection ladder]
For a declared ansatz \(\iota:\cA\hookrightarrow\cX_{\HS}\), a test map
\(P:\cY_{\HS}\to\cE\), and a flow \(\Phi_t\) when one is supplied, define:
\begin{enumerate}[label=\textnormal{L\arabic*:},leftmargin=3.4em]
\item \emph{projected consistency}: \(P\cR_{\HS}(\iota(a))=0\);
\item \emph{full equation closure}: \(\cR_{\HS}(\iota(a))=0\);
\item \emph{ansatz-local isolation}: \(a\) is isolated among full solutions
      in \(\cA\), modulo the declared equivalences;
\item \emph{dynamical attraction}: \(\iota(a)\) is an attracting fixed point
      of a specified evolution on a specified domain;
\item \emph{physical selection}: a source rule assigns the realized branch
      or a probability law on the relevant basins.
\end{enumerate}
\end{definition}

The ladder is not meant to devalue the lower levels.  A projected calculation
can reject an ansatz cheaply.  Full closure can construct a genuine solution.
Isolation can remove continuous moduli in the chosen family.  Attraction can
show robustness.  The point is that the noun ``selection'' must say which
level has been reached.

\subsection{Why component equations are not automatically complete}

\begin{proposition}[Projected residual criterion]\label{prop:projected}
For \(a\in\cA\),
\[
P\cR_{\HS}(\iota(a))=0
\quad\Longrightarrow\quad
\cR_{\HS}(\iota(a))\in\ker P.
\]
Consequently, projected consistency implies full equation closure if and only
if
\[
\im(\cR_{\HS}\circ\iota)\cap\ker P=\{0\}.
\tag{3.1}\label{eq:complete-test}
\]
\end{proposition}

\begin{proof}
The first statement is the definition of the kernel.  The implication to a
zero full residual holds precisely when the only residual both produced by
the ansatz and invisible to \(P\) is the zero residual.
\end{proof}

Condition \cref{eq:complete-test} is the algebraic form of a consistent
truncation check.  Expanding the Bianchi identity on a basis of invariant
\((2,2)\)-forms is complete for that residual only if the residual is known
to lie entirely in the tested invariant space.  It does not by itself prove
the Dolbeault integrability of a printed connection, the existence of global
bundle transition maps, stability, a global gerbe, or the discarded field
equations.

\begin{example}
Take \(\cY_{\HS}=\R^2\), let \(P(x,y)=x\), and let a proposed ansatz produce
\(\cR_{\HS}(\iota(a))=(0,1)\).  The tested coefficient vanishes, but the full
residual does not.  Heterotic analogues arise whenever one checks invariant
Bianchi coefficients while leaving a nonzero Maurer--Cartan, HYM, global
bundle, or non-invariant residual untested.
\end{example}

\subsection{Why isolated in an ansatz is not globally unique}

\begin{theorem}[Ansatz isolation theorem]\label{thm:ansatz}
Assume \(a_\star\in\cA\) satisfies the full equations and that the derivative
\[
D(\cR_{\HS}\circ\iota)_{a_\star}:T_{a_\star}\cA
\longrightarrow\cY_{\HS}
\]
is injective with a bounded left inverse on its image.  Then \(a_\star\) is
locally isolated in \(\cA\), modulo any quotient directions removed from
\(T_{a_\star}\cA\).  This conclusion does not imply local or global
uniqueness in \(\cX_{\HS}\).
\end{theorem}

\begin{proof}
The inverse or implicit function theorem gives local isolation on the
declared finite-dimensional slice after quotient directions are removed.
Directions transverse to \(\iota(\cA)\) are absent from the derivative.
Other full solutions may therefore lie arbitrarily near the slice, or on
disconnected components not meeting it at all.
\end{proof}

To promote ansatz uniqueness to global uniqueness, one would at minimum need
an exhaustion theorem showing that every admissible solution orbit meets the
ansatz and a faithfulness theorem showing that distinct full orbits are not
identified there.  Such statements are rare and are not consequences of
Nomizu reduction alone.

\subsection{What attraction and physical selection add}

\begin{definition}[Dynamical and physical selection data]
Dynamical attraction is typed by a semiflow
\(\Phi_t:\cD\to\cD\subseteq\cX_{\HS}\), a fixed point \(x_\star\), and a
basin
\[
\cB(x_\star)=\{x\in\cD:\Phi_t(x)\to x_\star\}.
\]
Physical selection additionally requires a source object, for example an
initial measure \(\mu_0\), a deterministic branch functional, or an
equivalent preparation law.  With a measure, the basin weight is
\(\mu_0(\cB(x_\star))\).
\end{definition}

\begin{theorem}[No selection from equations alone]\label{thm:no-selection}
The zero set \(\cS_{\HS}\) determines neither a semiflow, its basins, nor a
measure on those basins.  Even a singleton zero set inside \(\cA\) supplies
no physical-selection probability without additional source data.
\end{theorem}

\begin{proof}
Distinct vector fields can share the same zero set while assigning a fixed
point different stability types.  For example, on \(\R\), both
\(\dot x=-x\) and \(\dot x=x\) have zero set \(\{0\}\), but the first attracts
and the second repels.  Likewise, infinitely many mutually singular measures
can be placed on a fixed family of basins.  Neither object is determined by
the equation \(x=0\).
\end{proof}

\section{Where MTT enters}

MTT aims to supply structure above the lower equations: an admissible upper
carrier, a reduction map, an evolution, and fixed-point or continuation
rules.  That ambition is directly relevant to vacuum selection, but only
after the upper and lower objects are connected.

Let \(\cM\) be an MTT configuration space with vector field \(V_{\MTT}\), and
let
\[
\mathfrak{B}:\cM\longrightarrow\cX_{\HS}
\]
be a typed bridge.  If a lower Hull--Strominger evolution \(V_{\HS}\) is
declared, the essential compatibility equation is
\[
D\mathfrak{B}_m\,V_{\MTT}(m)
=V_{\HS}(\mathfrak{B}(m)).
\tag{4.1}\label{eq:intertwine}
\]
The adjacent MTT-to-Hull--Strominger paper owns the exact and residual
fixed-point descent theorems based on this equation
\cite{NeroMTTStrominger2026}.  We do not duplicate them here.  Their lesson
for selection is simple: an upper fixed point descends to a lower fixed point
only after the bridge and the flow intertwining are supplied.  A rank match,
shared notation, or formal similarity of constraints is not enough.

There are nevertheless narrower senses in which MTT already selects data.
The current finite program proves an exact \(q=79\) arithmetic branch under
its declared discrete assumptions, and it selects a retarded representative
within a conjugate orientation orbit.  Those are genuine finite branch
selections.  They do not by themselves select a smooth compactification,
because the visible and hidden bundles, common HYM chamber, differential
Bianchi representative, and continuum intertwiner remain separate objects.

\section{Re-auditing the two former case studies}

\subsection{Iwasawa: valid local geometry, invalid physical chain}

Let \(X=\Gamma\backslash H_3(\C)\) be the Iwasawa manifold with invariant
\((1,0)\)-forms
\[
\dd\omega^1=\dd\omega^2=0,\qquad
\dd\omega^3=\omega^1\wedge\omega^2.
\]
For a diagonal Hermitian form
\[
\omega=\frac{\mathrm{i}}{2}
\sum_{j=1}^{3}r_j^2\omega^j\wedge\overline{\omega^j},
\]
the local complex-parallelizable and balanced calculations are legitimate.
The torsion \(H=\mathrm{i}(\bar\partial-\partial)\omega\) and its derivative
can be computed exactly in the invariant frame.

The former paper then attached a rank-three visible-bundle construction,
component anomaly match, generation count, and normalized cubic Yukawa.
The current audited calculation shows why that chain fails
\cite{NeroFluxAudit2026}:
\begin{enumerate}[leftmargin=2em]
\item a form used as Chern data is not closed;
\item the advertised monad lacks globally constructed line bundles and
      holomorphic maps;
\item the printed Dolbeault matrix has a nonzero Maurer--Cartan residual;
\item the unique one-entry repair in the stated matrix ansatz lies in one
      complex-gauge orbit and has a non-scalar holomorphic commutant, so it
      does not supply the claimed stable simple rank-three bundle;
\item the stated trivial carrier forces \(c_3=0\);
\item the Bianchi, generation, and Yukawa conclusions depending on that
      bundle therefore do not follow.
\end{enumerate}

This is not a no-go theorem for all invariant Iwasawa solutions.  Such
solutions exist in the literature for carefully chosen connections and
bundles \cite{FeiYau2015,OtalUgarteVillacampa2016}.  It is a no-go for using
the particular invalid bundle chain as the selection witness in this paper.

\subsection{Lens--Nil: an auxiliary real structure}

The former second example used
\[
L(3,1)\times(\Gamma\backslash\mathrm{Nil}_3)
\]
with a balanced real \(\SU(3)\)-structure but explicitly found
\(\dd\Omega\neq0\).  The Hull--Strominger system requires a complex
threefold with integrable complex structure and holomorphic volume form.
Thus the displayed Lens--Nil object is not a point of \(\cX_{\HS}\), and
solving two formal coefficient equations on it cannot select a heterotic
vacuum.

Lens and Nil data may still be useful as auxiliary transport, filtration,
spectral, or finite-carrier models inside MTT.  That role is different from
identifying their literal product with the selected global compactification.
In particular, it is not the \(q=79\) Fu--Yau-oriented carrier.

\subsection{Status table}

\begin{center}
\small
\begin{tabular}{p{0.31\textwidth}p{0.18\textwidth}p{0.41\textwidth}}
\toprule
Object & Current status & Meaning for selection\\
\midrule
Diagonal Iwasawa geometry and local torsion & retained &
Valid lower geometric calculation; no physical bundle is selected.\\
Old Iwasawa visible bundle & withdrawn &
Cannot source anomaly, index, or Yukawa conclusions.\\
Old normalized Iwasawa Yukawa & withdrawn &
Normalization of \(\Omega\) does not determine normalized matter
wavefunctions or their physical overlap.\\
Lens--Nil compactification & retired as physical &
Non-integrable auxiliary real \(\SU(3)\)-structure only.\\
\(q=79\) arithmetic branch & exact at finite tier &
Selects a discrete branch under stated assumptions, not a smooth vacuum.\\
Rank-two Cech/HYM packets & finite or rank-two certified &
Useful ingredients; no automatic rank-three physical transfer.\\
Physical visible--hidden Fu--Yau tuple & open &
Required before compactification selection can be claimed.\\
\bottomrule
\end{tabular}
\end{center}

\section{A correct theorem for finite ansatz calculations}

The finite calculation can still support a rigorous theorem when its scope
is stated honestly.

\begin{theorem}[Finite-ansatz conclusion theorem]\label{thm:finite}
Let \(\cA\) be a finite-dimensional ansatz and let \(a_\star\in\cA\).
Assume:
\begin{enumerate}[label=\textnormal{(F\arabic*)},leftmargin=3.4em]
\item \(\iota(a)\) defines global typed data in \(\cX_{\HS}\) for every
      \(a\) in a neighborhood of \(a_\star\);
\item the tested residual map \(P\) satisfies the completeness condition
      \cref{eq:complete-test};
\item quantization, stability, gauge quotient, and global bundle or gerbe
      conditions are included rather than inferred from local forms;
\item \(P\cR_{\HS}(\iota(a_\star))=0\);
\item the quotient derivative at \(a_\star\) is injective with a bounded
      left inverse.
\end{enumerate}
Then \(\iota(a_\star)\) is a full Hull--Strominger solution locally isolated
inside the declared ansatz, modulo the declared equivalences.

If, in addition, a flow preserving \(\iota(\cA)\) is specified and its
linearization has a certified spectral gap with nonlinear remainder
controlled on a neighborhood, then \(a_\star\) is locally attracting in that
neighborhood.  Neither conclusion establishes global uniqueness or physical
selection without exhaustion and source hypotheses.
\end{theorem}

\begin{proof}
By (F2) and (F4), \cref{prop:projected} gives full equation closure.
Assumptions (F1) and (F3) ensure that this zero is a zero of the correctly
typed global problem rather than a local symbolic surrogate.  Assumption
(F5) and \cref{thm:ansatz} give isolation in the quotient ansatz.  The final
claim is the standard linearized-stability conclusion once the flow, gap,
domain invariance, and nonlinear bound are supplied.  The missing global and
source conclusions are excluded by \cref{thm:ansatz,thm:no-selection}.
\end{proof}

This theorem replaces the former ``FCC equals component equations''
statement.  It is stronger because its hypotheses reveal exactly what a
candidate must provide, and weaker in the necessary way because it does not
rename consistency as selection.

\section{The current q79 compactification boundary}

\subsection{What is already available}

The selected \(q=79\) program contains nontrivial and reusable ingredients:
\begin{itemize}[leftmargin=2em]
\item an exact arithmetic theorem selecting \(q=79\) on the declared finite
      branch;
\item a literal finite Cech witness;
\item a certified finite projected HYM approximation;
\item a selected rank-two continuum HYM witness with Fourier-tail and Wiener
      contraction control;
\item a retarded orientation representative inside the finite conjugate
      branch pair.
\end{itemize}
These facts improve the search problem.  They constrain which branch and
which local analytic architecture should be used.  They do not yet satisfy
the same-carrier requirement of \cref{eq:balanced,eq:hym,eq:bianchi}.

\subsection{Physical compactification-selection certificate}

A future claim that MTT selects one heterotic compactification should provide
at least the following certificate:
\begin{enumerate}[label=\textnormal{C\arabic*:},leftmargin=3.4em]
\item one global complex \(q=79\) Fu--Yau-oriented carrier with its balanced
      metric and holomorphic volume form;
\item explicit stable rank-three visible and compatible hidden holomorphic
      bundles in one positive Gauduchon/HYM chamber;
\item HYM connections and the differential Green--Schwarz identity on that
      same tuple, together with global flux/gerbe data;
\item a connection-preserving MTT bridge and the flow intertwining
      \cref{eq:intertwine}, or a certified residual theorem;
\item a basin or contraction proof on a declared domain if dynamical
      attraction is claimed;
\item an exhaustion theorem or an explicitly restricted candidate class if
      uniqueness is claimed;
\item a preparation, branch, or measure rule if physical realization or
      probabilities are claimed;
\item matter cohomology, normalized overlap data, and the worldsheet/IR
      endpoint for phenomenological claims.
\end{enumerate}

The first three rows are the open physical Hull--Strominger endpoint.  The
fourth is the open geometry-to-operator naturality problem.  Later rows
should not be used to hide those earlier dependencies.

\section{Worldsheet and effective-field-theory boundaries}

The first version claimed a succinct worldsheet verification, but the
displayed beta functions mixed bosonic-string central-charge notation with a
heterotic ten-dimensional target and treated the dilaton equation as
automatic.  That conclusion is withdrawn.  Supersymmetry, anomaly
cancellation, equations of motion, and worldsheet conformal invariance are
related but convention- and order-dependent statements.  The choice of
tangent connection also participates in field redefinitions and the anomaly
equation \cite{DelaOssaSvanes2014}.

For a valid Hull--Strominger tuple, leading-order sigma-model reasoning
provides an important consistency check.  It does not repair a nonintegrable
target or an undefined bundle, and it does not by itself prove an all-orders
SCFT.  Fu--Yau geometry and Anomaly-flow results provide established
mathematical routes to lower fixed points
\cite{FuYau2008,PhongPicardZhang2018}; they do not supply the missing MTT
source map automatically.

The former EFT statements are narrowed for the same reason.  A physical
Yukawa coupling requires identified matter cohomology classes, normalized
zero modes, bundle-valued products, a Hermitian inner product, and a
four-dimensional normalization and matching convention.  Rescaling
\(\Omega\) or choosing a basis in an invariant form space cannot set that
physical coupling to one.  Separate finite/profile Yukawa results elsewhere
in the MTT corpus are unaffected, but they cannot be retroactively attributed
to the invalid Iwasawa construction.

\section{Why the corrected result remains useful}

The correction changes the role of finite invariant calculations without
making them pointless.

\paragraph{They are efficient obstruction detectors.}
A nonzero projected residual disproves a candidate immediately.  Failure of
integrability, closure, stability, or a characteristic-class condition can
retire an entire branch before expensive continuum work.

\paragraph{They can produce candidate seeds.}
When the global data exist and the truncation is complete, a finite solution
can seed an implicit-function, perturbative, Galerkin, or flow construction.
The literature on stable bundles, Fu--Yau geometry, and invariant solutions
shows several ways in which such seeds can become genuine solutions.

\paragraph{They make the remaining assumptions visible.}
The selection ladder prevents a common source of apparent progress:
relabeling a solved constraint as an attractor or a finite branch as the
realized universe.  In MTT this is especially valuable because the theory is
explicitly attempting to move upstream from lower equations to their
carrier, evolution, and source.

\paragraph{They preserve narrower exact selections.}
The exact \(q=79\) arithmetic and retarded-orientation choices are not erased
because the physical compactification is open.  They are retained as
properly typed finite selections and can constrain the eventual global
construction.

\section{Conclusion}

The corrected conclusion is precise.  Componentwise left-invariant equations
can establish projected consistency.  With global typed data and a complete
residual test they can establish a full solution.  With a nonsingular
quotient derivative they can establish isolation inside the ansatz.  With a
specified flow and certified stability bounds they can establish attraction.
Only an exhaustion statement and a source or preparation rule can support a
claim of physical vacuum selection.

The former Iwasawa and Lens--Nil examples do not reach those levels: the
Iwasawa physical bundle chain fails, and Lens--Nil is not a complex
Hull--Strominger target.  The exact \(q=79\) finite branch and rank-two
Cech/HYM results remain real progress at narrower tiers.  The next physical
advance is therefore not another reformulation of the component equations.
It is the common rank-three visible--hidden Fu--Yau tuple, followed by the
connection-preserving MTT bridge, its attraction domain, and its source rule.



\bibliographystyle{plain}

% BEGIN MTT MANAGED COMPUTATIONAL EVIDENCE
\section*{Computational Evidence and Reproducibility}
The q79 theorem and audit, literal finite rank-two Cech witness, and rank-two Wiener-contraction certificate are direct evidence only for the finite branch, topological witness, and declared rank-two analytic tier. They do not establish global exhaustion, a physical rank-three visible-hidden Hull-Strominger tuple, dynamical attraction of that tuple, or a source measure selecting one realized compactification.

The referenced rows are frozen to the curated results repository state identified below. Hashes are grouped in eight-character blocks for line breaking.
\begin{quote}\small
\textbf{Repository:} \url{https://github.com/PeterNero/mtt-results-repro}\\
\textbf{Commit:} \texttt{31247ebb 5c22f3fb b5443024 365433c6 ee0bff4a}\\
\textbf{Manifest:} \path{release/result_manifest.json}\\
\textbf{Manifest SHA-256:}\\
\texttt{fb399689 60b00584 631dbf53 1a708e18}\\
\texttt{ef928d6b 6d935119 c185d7f6 32b1e7cd}
\end{quote}

Tier labels are quoted verbatim from that manifest. A row used directly supports only the specific computational statement identified above; a corpus-state cross-check does not prove this paper's local theorems; and an open row is evidence of an unresolved obligation, never of closure.

\par\begingroup\dimen0=\pagegoal\advance\dimen0 by -\pagetotal\ifdim\dimen0<7\baselineskip\newpage\fi\endgroup
\paragraph{Rows used directly in this paper.}
\begin{itemize}
\setlength{\itemsep}{0.25em}
\item \path{A19/hym_wiener_contraction} (\emph{numeric certified}).\par Weighted-theta Fourier-tail and Wiener contraction certificate.
\item \path{A07/literal_cech_witness} (\emph{derived exact}).\par Literal 81-entry, 729-cocycle finite Cech witness.
\item \path{A11/q79_exact_audit} (\emph{derived exact}).\par Executable q=79 exact-branch audit.
\item \path{A11/q79_exact_theorem} (\emph{derived exact}).\par CRT q=79 theorem on the selected exact branch.
\end{itemize}

No imported row changes theorem ownership or promotes a neighboring claim: all local statements retain their stated hypotheses, domains, and limitations.
% END MTT MANAGED COMPUTATIONAL EVIDENCE

\bibliography{main}

\end{document}
